\renewcommand{\theposttitle}{How to use evidence to build confidence}
\renewcommand{\thepostdate}{2025-12-30}
\vspace*{\fill}
\begin{center}
\includegraphics[width=\textwidth, height=0.55\textheight, keepaspectratio]{C:/Users/James/Documents/GitHub/VisualIdeasBook/data/images/e83c09ee515f52266cffce3137af6ee0.png}
\end{center}
\vspace{2mm}
\begin{center}
\begin{minipage}{0.92\textwidth}
\small
For years, runners believed the human body couldn't run a mile in under four minutes. Doctors said it was physically impossible!

Then on May 6, 1954, British runner Roger Bannister ran it in 3:59.4.

But here’s the crazy part: 46 days later, John Landy from Australia did it in 3:58.0. By year's end, three more runners beat the “impossible” 4-minute mile.

Nothing changed physically, except belief.

I love this story because whenever someone does the "impossible,” others see themselves in that someone.

“I can't" → "they did"

"They did" → "I might"

"I might" → "I will"

It's like a \textbf{possibility chain of evidence}. Of course, your "four-minute mile" might look different. Maybe it's writing that post and hitting publish. Or speaking up in meetings. Or just facing the day.

And it works both ways. You might need someone else's proof to believe in what you can do. Someone else needs your proof to believe in what they can do.

Either way, you're part of this chain.

Break your barriers. Not just for you. \\
For everyone watching who thinks they can't.

(if you like this idea, check out my visual for The Stockdale Paradox)
\end{minipage}
\end{center}
\vspace*{\fill}